% Part: first-order-logic
% Chapter: beyond
% Section: many-sorted-logic

\documentclass[../../../include/open-logic-section]{subfiles}

\begin{document}

\olfileid{fol}{byd}{msl}

\olsection{کئی قسماں والی منطق}

پہلے درجے دی منطق وچ متغیر تے مسوّر اک ہی
!!{domain} دے عناصر اُتے چلدے نیں۔ پر اکثر کئی
وکھ وکھ، اک دوجے توں جدا !!{domain}s رکھنا فائدہ مند
ہوندا اے: مثال لئی اعداد دا !!a{domain}،
ہندسی شےواں دا !!a{domain}، اعداد توں اعداد ول جان
والے فنکشناں دا !!a{domain}، ابیلی گروپاں دا
!!a{domain}، تے ایس طرح ہور۔

کئی قسماں والی منطق ایس کم لئی اک ڈھانچا دیندی اے۔
شروع وچ ”قسماں“ (sorts) دی فہرست لئی جاندی اے: کسے شے
دی ”قسم“ (sort) ایہ دسدی اے کہ اوہ کیہڑے
”!!{domain}“ وچ آندی اے۔ فیر ہر قسم لئی وکھرے
!!{variable}s تے مسوّر ہوندے نیں، تے (عموماً) ہر
قسم لئی اپنی !!{identity} وی ہوندی اے۔ فنکشناں تے
تعلقاں دے دلیل والے مقام وی قسماں دے حساب نال
مقرر ہوندے نیں، یعنی ہر مقام اُتے کیہڑی قسم دی شے
آ سکدی اے۔ ایس توں علاوہ پہلے درجے دی منطق دے
معمول دے قاعدے برقرار رہندے نیں؛ مسوّراں دے
قاعدیاں دی ہر قسم لئی اپنی صورت دہرائی جاندی اے۔

مثال لئی بین الاقوامی تعلقاں دے مطالعے واسطے اسیں
ایسی زبان چن سکدے آں جیہدیاں شےواں دیاں دو قسماں
ہون: فرانسیسی شہری تے جرمن شہری۔ پہلی قسم دیاں
شےواں لئی اک اک-جگہی تعلق ”انگوری شراب پیندا اے“ رکھ
سکدے آں؛ دوجی قسم دیاں شےواں لئی دوجا اک-جگہی
تعلق ”ورسٹ کھاندا اے“؛ تے دو-جگہی تعلق ”دو
ملکاں توں بنیا ویاہتا جوڑا بناندے نیں“، جیہدے
دو دلیل نیں: پہلی دلیل پہلی قسم دی، تے دوجی
دوجی قسم دی۔ جے متغیر $a$, $b$, $c$ فرانسیسی
شہریاں اُتے تے $x$, $y$, $z$ جرمن شہریاں اُتے
چلدے نیں، تاں
% PNB-SYN-009: اندرلے مسوّر دے متغیر تے ماتحت فارمولے نوں اوہناں دے اپنے اختیاری آرگومنٹاں وچ رکھیا۔
\[
\lforall[a][\lforall[x][(\Atom{\Obj{MarriedTo}}{a,x} \lif
(\Atom{\Obj{DrinksWine}}{a} \lor \lnot \Atom{\Obj{EatsWurst}}{x})]]
\]
ایہ بیان کردا اے کہ جے کوئی فرانسیسی شخص کسے
جرمن نال ویاہیا اے، تاں یا اوہ فرانسیسی شخص
انگوری شراب پیندا اے یا اوہ جرمن شخص ورسٹ نہیں کھاندا۔

کئی قسماں والی منطق نوں قدرتی طریقے نال پہلے درجے
دی منطق وچ سمویا جا سکدا اے: کئی قسماں دے
!!{domain}s دیاں ساریاں شےواں نوں اک پہلے درجے دے
!!{domain} وچ اکٹھا کر دیو، قسماں دی پہچان رکھن
لئی اک-جگہی !!{predicate}s ورتو، تے مسوّراں دا
دائرہ متعلقہ قسم تک محدود کرو۔ مثال لئی اوپرلی
مثال نال ملدی پہلے درجے دی زبان وچ پہلے بیان
کیتے تعلقاں دے نال اک-جگہی !!{predicate}s
”$\Obj{German}$“ تے ”$\Obj{French}$“ وی ہون گے،
پر تعلقاں دی قسم-بندی دیاں شرطاں علامتاں توں
ہٹا دتیاں جان گیاں۔ اک قسم-بند مسوّر
$\lforall[x][!A]$ دا ترجمہ، جتھے $x$ جرمن
قسم دا !!a{variable} اے، ایہ بنے گا:
\[
\lforall[x][(\Atom{\Obj{German}}{x} \lif !A)].
\]
اسیں ایداں دے مسلّمے شامل کرن ہون گے جیہڑے
قسماں دے اک دوجے توں جدا ہون دی ضمانت دین،
مثلاً $\lforall[x][\lnot (\Atom{\Obj{German}}{x} \land
  \Atom{\Obj{French}}{x})]$؛ نال ای ایداں دے
مسلّمے وی جیہڑے ضمانت دین کہ ”انگوری شراب پیندا اے“
صرف انہاں شےواں لئی درست ہووے جیہڑیاں
$\Atom{\Obj{French}}{x}$ محمول پورا کردیاں نیں،
وغیرہ۔
% PNB-SEM-028: معمول دے غیر خالی قسمی دائرے، ایس دو-قسمی مثال دی مکمل ڈھکائی، تے مستقل/فنکشن دی قسم-بندی واضح کیتی۔
ایتھے معمول دے قاعدیاں لئی ہر قسم دا دائرہ غیر خالی
منیا گیا اے۔ ایس دو-قسمی مثال وچ اکٹھے دائرے دا ہر
عنصر فرانسیسی یا جرمن قسم وچوں کسے اک دا ہونا چاہیدا
اے۔ جے زبان وچ مستقل یا فنکشن ہون، تاں مستقل اپنی
مقرر قسم وچ ہوون، تے فنکشن مقرر انپٹ قسماں دے عناصر
نوں اپنی مقرر آؤٹ پٹ قسم دا عنصر دین۔ فنکشناں دیاں
مقرر انپٹ قسماں توں باہر دیاں قدراں ایس ترجمے وچ
بےاثر نیں۔ انہاں رواجاں تے مسلّمیاں نال ایہ وکھانا
اوکھا نہیں کہ کئی قسماں والے !!{sentence}s دا
ترجمہ پہلے درجے دے !!{sentence}s وچ ہوندا اے،
تے کئی قسماں والے !!{derivation}s دا ترجمہ
پہلے درجے دے !!{derivation}s وچ۔ ایس طرح کئی
قسماں والیاں !!{structure}s دا وی متعلقہ پہلے
درجے دیاں !!{structure}s وچ ”ترجمہ“ ہوندا اے،
تے الٹ پاسے وی؛ ایس لئی کئی قسماں والی منطق
لئی تمامیت دا قضیہ وی مل جاندا اے۔

\end{document}
